% xelatex compte-rendu-de-TP
% modifier le document à partir de la ligne 90

\documentclass[11pt]{article}
\DeclareTextCommandDefault{\nobreakspace}{\leavevmode\nobreak\ }
\usepackage{xunicode}
\usepackage{fontspec}
\usepackage{xltxtra}
\usepackage{url}
% \usepackage[francais]{babel}
\usepackage{fullpage}
\usepackage{fancyvrb}
\usepackage{tikz}
\usepackage{xcolor}
\usepackage{listings}

\definecolor{mygreen}{rgb}{0,0.6,0}
\definecolor{mygray}{rgb}{0.5,0.5,0.5}
\definecolor{mymauve}{rgb}{0.58,0,0.82}
\definecolor{terminalbgcolor}{HTML}{330033}
\definecolor{terminalrulecolor}{HTML}{000099}
\newcommand{\lstconsolestyle}{
\lstset{
	backgroundcolor=\color{terminalbgcolor},
	basicstyle=\color{white}\fontfamily{lmtt}\footnotesize\selectfont,
	breakatwhitespace=false,  
	breaklines=true,
	captionpos=b,
	commentstyle=\color{mygreen},
	deletekeywords={...},
	escapeinside={\%*}{*)},
	extendedchars=true,
	frame=single,
	keepspaces=true,
	keywordstyle=\color{blue},
	%language=none,
	morekeywords={*,...},
	numbers=none,
	numbersep=5pt,
    framerule=2pt,
	numberstyle=\color{mygray}\tiny\selectfont,
	rulecolor=\color{terminalrulecolor},
	showspaces=false,
	showstringspaces=false,
	showtabs=false,
	stepnumber=2,
	stringstyle=\color{mymauve},
	tabsize=2
}
}

\lstconsolestyle

%% Environnements

%% Ensembles mathématiques

\newcommand{\A}{\mathbb{A}}
\newcommand{\Q}{\mathbb{Q}}
\newcommand{\Z}{\mathbb{Z}}
\newcommand{\C}{\mathbb{C}}
\newcommand{\N}{\mathbb{N}}
\newcommand{\R}{\mathbb{R}}
\renewcommand{\emptyset}{\varnothing}

\makeatletter %
\renewcommand{\verbatim@font}{ %
    \ttfamily\footnotesize\catcode`\<=\active\catcode`\>=\active %
}
\makeatother

\begin{document}
\begin{center}
\textbf{\bf Rapport de Projet d'ALN}
 
\end{center}
\medskip
\hrule
 
\bigskip

\section{Résolution de systèmes triangulaires}

J'ai programmé et testé ma propre version 
(nommée \texttt{my\_solve\_triangular}) de \texttt{solve\_triangular}. 
Son paramétrage est compatible avec celui de \texttt{solve\_triangular}.
Note~: je n'ai pas eu le temps de tout programmer.
L'exception \texttt{ValueError ("pas encore implémenté")} est
levée si une fonctionnalité non implémentée est appelée.

Fichier~: \texttt{projet/mini-exemple/my\_solve\_triangular.py}. 

Fichier de test~: \texttt{my\_solve\_triangular.py}.

Extrait de la documentation~:
\begin{Verbatim}[fontsize=\small]
# Résout A . x = b où A est une matrice triangulaire
#         trans : 'N' ou 'T' (indique s'il faut transposer A ou pas)
#         lower : True ou False (indique où sont les données)
# unit_diagonal : True ou False
# Retourne le vecteur x solution

def my_solve_triangular (A, b, trans='N', lower=True, unit_diagonal=False) 
\end{Verbatim}

Exécution des tests~:
\begin{lstlisting}
$ python3 test_my_solve_triangular.py 
Exécution de test_my_solve_triangular_1 : succès
Nombre de tests  : 1
Nombre de succès : 1
Nombre d'échecs  : 0
\end{lstlisting}

\end{document}


